\documentclass[11pt,a4paper]{article}

\usepackage[utf8]{inputenc}
\usepackage[T1]{fontenc}
\usepackage[french]{babel}
\usepackage[a4paper,margin=2.4cm]{geometry}
\usepackage{lmodern}
\usepackage{xcolor}
\usepackage{graphicx}
\usepackage{enumitem}
\usepackage{booktabs}
\usepackage{array}
\usepackage{titlesec}
\usepackage{fancyhdr}
\usepackage[hidelinks]{hyperref}

% Couleurs (orange de l'application)
\definecolor{accent}{HTML}{D97757}
\definecolor{dark}{HTML}{243857}
\definecolor{light}{HTML}{F5F3EE}

\hypersetup{colorlinks=true, linkcolor=dark, urlcolor=accent}

% Titres de section colorés
\titleformat{\section}{\Large\bfseries\color{dark}}{\thesection}{0.6em}{}
\titleformat{\subsection}{\large\bfseries\color{accent}}{\thesubsection}{0.5em}{}

% En-tête / pied de page
\pagestyle{fancy}
\fancyhf{}
\renewcommand{\headrulewidth}{0.4pt}
\fancyhead[L]{\small\color{dark}Voyage Assistant}
\fancyhead[R]{\small\color{dark}SAE2 BUT3}
\fancyfoot[C]{\thepage}

\newcommand{\code}[1]{\texttt{\small #1}}

\begin{document}

% ----------------------------------------------------------------------
% Page de garde
% ----------------------------------------------------------------------
\begin{titlepage}
\centering
\vspace*{1.5cm}
{\color{accent}\rule{\linewidth}{2pt}}\\[0.8cm]
{\Huge\bfseries\color{dark} Voyage Assistant}\\[0.4cm]
{\Large Assistant IA conversationnel pour le transport}\\[0.3cm]
{\large Rapport de projet --- SAE2 BUT3}\\[0.5cm]
{\color{accent}\rule{\linewidth}{2pt}}\\[1.2cm]

\begin{tabular}{>{\bfseries}l l}
Ako Christian & IA conversationnelle \& RAG, déploiement \\[0.2cm]
Carl Similien & Back-end \& intégration (API, base de données) \\[0.2cm]
Noé Cervera & Modèles \& évaluation (embeddings, reranker, NER) \\[0.2cm]
Dhanoush Kessavane & Front-end, design \& documentation \\
\end{tabular}\\[1.2cm]

{\large Encadrement : \textbf{Bilal Faye} (LIPN, CNRS UMR 7030)}\\[0.6cm]
{\large Promotion 2025--2026}\\[1.5cm]

\fbox{\parbox{0.9\linewidth}{\centering\vspace{0.2cm}
\textbf{Application en ligne}\\[0.2cm]
\url{https://optimalako-voyage-assistant.hf.space}\\[0.2cm]
\textbf{Code source}~: \url{https://github.com/optmlako2004/chatbot}
\vspace{0.2cm}}}

\vfill
\end{titlepage}

\tableofcontents
\newpage

% ----------------------------------------------------------------------
\section{Introduction et contexte}

Ce projet s'inscrit dans la \textbf{SAE2 du BUT3}, dont l'objectif est de concevoir
une application complète intégrant un \textbf{assistant conversationnel} fondé sur un
modèle de langage (LLM) et une approche \textbf{RAG} (\emph{Retrieval-Augmented
Generation}), spécialisée sur un domaine métier précis. Nous avons choisi le domaine
du \textbf{transport de voyageurs} (avion, train, bateau, bus longue distance).

L'application, baptisée \textbf{Voyage Assistant}, permet de rechercher et réserver
un trajet, de gérer son compte et ses billets, et surtout de \textbf{dialoguer avec
un assistant IA} qui répond aux questions de voyage et exécute réellement des actions
(modification de réservation, réclamation, annulation). Elle est \textbf{déployée en
ligne} et accessible publiquement.

\subsection{Objectifs}
\begin{itemize}[nosep]
  \item Construire une application \textbf{full-stack} (frontend, backend, base de données).
  \item Intégrer un \textbf{LLM} et une chaîne \textbf{RAG vectorielle} conforme aux concepts vus en cours.
  \item Garantir des réponses \textbf{factuelles} (anti-hallucination) en s'appuyant sur des sources de vérité.
  \item \textbf{Déployer} l'ensemble en ligne, gratuitement.
\end{itemize}

% ----------------------------------------------------------------------
\section{Cahier des charges et réponses apportées}

Le tableau suivant met en regard chaque exigence du sujet et sa réalisation concrète.

\begin{center}
\renewcommand{\arraystretch}{1.3}
\begin{tabular}{>{\bfseries}p{4.8cm} p{9.5cm}}
\toprule
Exigence & Réalisation \\
\midrule
Domaine spécialisé & Transport de voyageurs (avion / train / bateau / bus) \\
Application full-stack & FastAPI (back) + React (front) + PostgreSQL \\
LLM & Google Gemini 2.5 Flash \\
RAG / vectorisation & LangChain + FAISS + embeddings \code{gte-small} + reranker CrossEncoder \\
Base de session & Tables \code{users}, \code{chat\_sessions}, \code{chat\_messages} \\
Voix (option) & Web Speech API : dictée vocale + lecture des réponses \\
Recherche web (bonus) & DuckDuckGo + APIs temps réel (météo, heure, change) \\
Déploiement en ligne & Hugging Face Spaces + Neon (PostgreSQL) \\
\bottomrule
\end{tabular}
\end{center}

% ----------------------------------------------------------------------
\section{Architecture technique}

L'application suit une architecture \textbf{client--serveur} classique, enrichie d'une
couche d'intelligence artificielle côté serveur.

\subsection{Vue d'ensemble}
\begin{itemize}[nosep]
  \item \textbf{Frontend} (\code{frontend/}) : React 18 (sans étape de build, via Babel
        navigateur), CSS personnalisé avec thèmes clair/sombre, et \textbf{Web Speech
        API} pour la voix.
  \item \textbf{Backend} (\code{backend/}) : FastAPI, validation Pydantic, ORM
        SQLAlchemy, et les services d'IA (LLM, RAG, reranker, NER, APIs temps réel).
  \item \textbf{Base de données} : PostgreSQL (Neon) en production, SQLite en
        développement. Le code est agnostique grâce à une simple variable
        d'environnement \code{DATABASE\_URL}.
\end{itemize}

\subsection{Modèle de données}
Les entités principales sont : \code{User}, \code{Route} (catalogue de plus de 182\,000
liaisons réelles), \code{Trajet} (occurrence datée d'une route), \code{Billet},
\code{Reclamation}, ainsi que \code{ChatSession} et \code{ChatMessage} pour la
\textbf{persistance des conversations}.

% ----------------------------------------------------------------------
\section{L'intelligence conversationnelle}

Le c\oe ur du projet est l'assistant. Sa philosophie : \textbf{un LLM ne suffit pas}.
Seul, un modèle de langage \emph{invente} (hallucinations) et ignore les données
récentes ou personnelles. Nous l'avons donc couplé à plusieurs sources de vérité via
le \textbf{RAG}.

\subsection{LLM --- Google Gemini 2.5 Flash}
Le modèle reçoit un \emph{system prompt} strict (rôle, périmètre, refus codifiés),
l'historique de la conversation, le contexte récupéré, puis le message courant. Il ne
sert qu'à \textbf{rédiger} la réponse, jamais à exécuter une action sensible.

\subsection{Qu'est-ce que le RAG et la vectorisation ?}
\textbf{RAG} (\emph{Retrieval-Augmented Generation}) consiste à \textbf{récupérer
l'information pertinente avant de générer} la réponse, puis à l'injecter dans le
prompt. Pour retrouver cette information par le \emph{sens} et non par mots-clés, on
\textbf{vectorise} chaque texte : il est converti en un vecteur de nombres
(384 dimensions) qui encode sa signification. Deux phrases de sens proche ont des
vecteurs proches, \emph{même si elles n'emploient pas les mêmes mots}. C'est ce
mécanisme qui permet de relier « je veux annuler » à un article de CGV intitulé
« conditions de résiliation ».

\subsection{La chaîne RAG en deux étapes}
\begin{enumerate}[nosep]
  \item \textbf{Indexation} : les documents (CGV) sont découpés en \emph{chunks}
        (\code{RecursiveCharacterTextSplitter}), vectorisés par le modèle
        \code{thenlper/gte-small} (exécuté \textbf{en local}), et stockés dans un index
        \textbf{FAISS}.
  \item \textbf{Recherche} : pour chaque question, un \textbf{bi-encoder} remonte
        rapidement 10 candidats, puis un \textbf{cross-encoder}
        (\code{ms-marco-MiniLM-L-6-v2}) les \textbf{re-trie} finement (\emph{reranking})
        pour ne garder que les 3 plus pertinents.
\end{enumerate}
Cette double étape offre la précision du cross-encoder sans payer son coût sur tout le
corpus.

\subsection{Extraction d'entités (NER)}
Pour fluidifier les parcours sensibles, un modèle \textbf{NER}
(\code{Jean-Baptiste/camembert-ner}) extrait le prénom et le nom d'une phrase libre.
L'utilisateur peut ainsi tout dire d'un coup au lieu de répondre à plusieurs questions.

\subsection{Les trois sources de connaissance}
Selon la question, l'assistant interroge la bonne source :
\begin{itemize}[nosep]
  \item \textbf{RAG documentaire} (FAISS) pour les règles (annulation, bagages\dots).
  \item \textbf{Base de données} pour les questions personnelles. Exemple :
        « combien de temps va durer mon vol~? » $\rightarrow$ lecture du billet de
        l'utilisateur, \textbf{calcul de la durée} à partir des horaires, réponse
        personnalisée.
  \item \textbf{APIs temps réel et web} pour les données du jour : \textbf{Open-Meteo}
        (météo, heure locale, fuseau), \textbf{Frankfurter} (taux de change),
        \textbf{DuckDuckGo} (compléments). Les informations stables (capitale, monnaie
        d'un pays) ne sont pas appelées : le LLM les connaît déjà.
\end{itemize}

\subsection{Machine à états pour les actions réelles}
Modifier, annuler ou réclamer ne sont \textbf{jamais} confiés au LLM. Une
\textbf{machine à états} déterministe pilote le parcours : intention détectée
$\rightarrow$ vérification d'identité (numéro de billet, nom, prénom, date de
naissance) $\rightarrow$ exécution réelle en base $\rightarrow$ envoi de l'e-mail et du
billet PDF mis à jour. Cette séparation \textbf{exécution / rédaction} élimine les
hallucinations d'action.

% ----------------------------------------------------------------------
\section{Fonctionnalités principales}
\begin{itemize}[nosep]
  \item Recherche multimodale (aller simple / aller-retour) avec gestion des escales
        sur plus de 182\,000 liaisons.
  \item Réservation avec choix de classe et bagages, génération d'un \textbf{billet
        PDF} (ReportLab) et envoi d'un \textbf{e-mail de confirmation} (Brevo).
  \item Authentification \textbf{Google} ou e-mail / mot de passe (tokens signés).
  \item Assistant IA : questions, modification, réclamation, annulation, suivi.
  \item \textbf{Voix} : dictée et lecture des réponses (Web Speech API).
  \item Persistance et reprise des conversations, notation en fin de session.
\end{itemize}

% ----------------------------------------------------------------------
\section{Sécurité}
La sécurité du chatbot repose sur plusieurs couches : nettoyage des entrées, garde-fous
contre insultes et \textbf{injections de prompt}, anti-extraction d'informations sans
identité vérifiée, \emph{system prompt} restrictif, requêtes SQL paramétrées
(anti-injection), échappement HTML natif de React (anti-XSS), et triple vérification
d'identité. Aucun secret n'est présent dans le code : tout passe par des variables
d'environnement.

% ----------------------------------------------------------------------
\section{Déploiement}

\subsection{Choix d'hébergement}
La chaîne RAG impose \code{torch} et deux modèles Transformers en mémoire
(environ 700~Mo à 1~Go de RAM). Les offres gratuites de type Render ou Railway
(512~Mo) provoquent un dépassement mémoire. Nous avons donc retenu \textbf{Hugging Face
Spaces} (16~Go de RAM gratuits, conçu pour les modèles ML).

\subsection{Architecture de production}
Un \textbf{unique service} (conteneur Docker) sert à la fois l'API, le RAG et le
frontend statique --- même origine, donc aucun problème de CORS. La base de données est
externe : \textbf{Neon} (PostgreSQL serverless), volontairement co-localisée avec le
Space (région \emph{us-east-1}) pour minimiser la latence.

\subsection{Étapes clés}
\begin{itemize}[nosep]
  \item \textbf{Image Docker} : installation de \code{torch} en version \textbf{CPU}
        (200~Mo au lieu de 1,2~Go en CUDA), pré-téléchargement des modèles RAG, écoute
        sur le port 7860.
  \item \textbf{Migration des données} : les 182\,511 routes n'existaient que dans le
        fichier SQLite local ; un script dédié (\code{migrate\_to\_postgres.py}) les
        copie vers Neon.
  \item \textbf{Secrets} configurés côté Hugging Face : \code{DATABASE\_URL},
        \code{AUTH\_SECRET}, \code{GEMINI\_API\_KEY}, etc.
  \item \textbf{OAuth Google} : ajout de l'URL du Space dans les origines autorisées.
\end{itemize}

% ----------------------------------------------------------------------
\section{Difficultés rencontrées et solutions}

\begin{center}
\renewcommand{\arraystretch}{1.3}
\begin{tabular}{p{6.2cm} p{8.1cm}}
\toprule
\textbf{Difficulté} & \textbf{Solution} \\
\midrule
DuckDuckGo limité (\code{202 Ratelimit}) & Passage à la librairie \code{ddgs}
(rotation multi-backend automatique). \\
Google ne fournit pas la date de naissance & Le formulaire de réservation devient la
source de vérité ; vérification tolérante à l'inversion nom/prénom. \\
Hallucinations sur les actions & Séparation stricte machine à états / LLM. \\
Recherche d'escales lente (15~s) & Index SQL au démarrage + cache en mémoire
($\rightarrow$ 0,34~s). \\
RAG « classique » (retour encadrant) & Réécriture en LangChain + FAISS + embeddings
locaux + reranking + NER. \\
\code{torch} CUDA trop lourd (1,2~Go) & Installation de \code{torch} CPU dans le Docker. \\
512~Mo de RAM insuffisants & Choix de Hugging Face Spaces (16~Go). \\
Migration des 182\,k routes & Script de copie SQLite $\rightarrow$ Neon par lots. \\
\bottomrule
\end{tabular}
\end{center}

% ----------------------------------------------------------------------
\section{Répartition du travail}

\begin{center}
\renewcommand{\arraystretch}{1.3}
\begin{tabular}{>{\bfseries}l p{9.5cm}}
\toprule
Membre & Contributions \\
\midrule
Ako Christian & Conception de l'IA conversationnelle (LLM, chaîne RAG, reranker, NER,
APIs temps réel), développement backend et mise en production. \\
Carl Similien & Développement back-end (API FastAPI, base de données) et travail
d'intégration entre les modules. \\
Noé Cervera & Mise au point et évaluation des modèles (embeddings, reranking, NER),
qualité des réponses. \\
Dhanoush Kessavane & Développement du front-end, conception graphique et rédaction de
la documentation. \\
\bottomrule
\end{tabular}
\end{center}

% ----------------------------------------------------------------------
\section{Conclusion et perspectives}

Voyage Assistant répond à l'ensemble des exigences de la SAE2 : application full-stack,
LLM couplé à une chaîne RAG vectorielle conforme aux concepts du cours, persistance des
sessions, option vocale, recherche web enrichie d'APIs temps réel, et \textbf{déploiement
en ligne effectif}. Au-delà de la simple intégration d'un modèle, le projet démontre
une architecture où l'IA s'appuie sur des \textbf{sources de vérité} pour rester
factuelle et utile.

\textbf{Perspectives} : enrichir le RAG avec davantage de documents métier, ajouter des
notifications en temps réel, et industrialiser le déploiement (intégration continue,
montée en charge).

\end{document}
